\newcommand{\CaseID}{august\_complex\_2020}
\newcommand{\CaseTitle}{2020 August Complex fire reconstruction against GOFER observations}
\newcommand{\FireName}{August Complex}
\newcommand{\FireYear}{2020}
\newcommand{\EventStart}{16 August 2020 at 21:00~UTC}
\newcommand{\EventEnd}{16 October 2020 at 22:00~UTC}
\newcommand{\DurationHours}{1465}
\newcommand{\CRSName}{UTM zone 10N (EPSG:32610)}
\newcommand{\FuelVersionText}{LANDFIRE release 2.1.0}
\newcommand{\FireSummary}{The supplied GOFER event record contains 1465 hourly perimeters, a final area of 4342.79~km$^2$ and an official size of 1,032,648~ac. The configuration uses 12 timed ignition point(s).}
\newcommand{\ConfigurationText}{The supplied simulation runs for 1465~h, with an initial timestep of 1~s, target Courant--Friedrichs--Lewy number of 0.2, 1466 hourly weather fields and 12 ignition point(s). Arrival time, spread rate, fireline intensity and fire-size statistics are requested at the supplied dump interval. Crown-fire model 1 and firebrand spotting are enabled. Unspecified settings retain the inspected source defaults, including fixed random seed 2024, foliar moisture 90\% and maximum timestep 600~s. A single simulation member is configured.}
\newcommand{\WeatherOffsetText}{Simulation and weather time zero is 0~h after the GOES ignition hour; comparisons apply this offset explicitly.}
\newcommand{\ElmfireRevision}{8981f14e0170}
\newcommand{\ElmfireBranch}{yiren}
